\chapter*{Plain-language dictionary and symbol guide}
\addcontentsline{toc}{chapter}{Plain-language dictionary and symbol guide}
\markboth{Dictionary and symbol guide}{Dictionary and symbol guide}

\textbf{Action.} A declared physical transformation with a specified boundary, quantity or other input, predecessor state and successor. It can be a transfer or another represented change. Naming a useful activity does not by itself define its state effect.

\textbf{Actor.} The person, organization or modeled decision-maker associated with an action. Its identity is distinct from the action's field value. The field equation does not assign moral worth to the actor or recipient.

\textbf{Affected factors.} Terms in a potential that depend on at least one changed coordinate. A complete local calculation includes every such term. In the separable Gaussian model, these are the changed endpoint terms.

\textbf{Baseline, $z$.} The common state from which an action and its no-action alternative are compared. In the prospective forcing-first order it is the post-forcing, pre-action state. A baseline must identify both data and model version.

\textbf{Capacity account, $B_i$.} A proposed persistent account controlling later modeled affordability. It is not the physical stock, storage capacity, household income or an automatically legitimate entitlement. Its update and use require a separately declared rule.

\textbf{Cell or node.} An indexed location or subsystem holding a represented state. A cell is a modeling boundary rather than necessarily a household, person or geographic square.

\textbf{Common path, $\gamma(s)$.} A single declared state path used to calculate the group integral and its child contributions. The path parameter is not automatically clock time. A physical-path reading needs a physical model that realizes it.

\textbf{Conservation.} A complete physical quantity balance within a named boundary. Lossless internal transfers preserve the sum of homogeneous stocks. Transformations, sinks and boundary crossings must be named when they are present.

\textbf{Curvature.} The way a slope changes as the state changes. For this quadratic it is represented by the fixed matrix $H$. Curvature explains why an initial slope times a finite quantity differs from the exact finite value.

\textbf{EBU, $E$.} The signed before/after potential difference minus declared nonduplicated process burden. In the present model its field component is dimensionless standardized deviation reduction. The name does not make it physical energy or money.

\textbf{Edge.} A declared possible directed physical process. For a lossless transfer, its incidence column has $-1$ at the source and $+1$ at the destination. A line drawn between places is not sufficient evidence that a real process is available.

\textbf{Equilibrium.} A property of a specified dynamical system, typically a state unchanged by its evolution rule. It is distinct from a declared reference or a minimum of a chosen potential unless the required relationship is established.

\textbf{External audit, $J$.} The proposed cumulative signed potential change from external forcing. It explains externally introduced or removed deviation. It is not a spendable ``nature wallet.''

\textbf{Factor.} One term in a decomposed potential, with a declared set of input coordinates. A factor can be local without depending on only one coordinate. Its dependencies determine the required evaluation neighbourhood.

\textbf{Feasibility.} Satisfaction of the declared physical conditions on an action, such as nonnegative stock, permitted edges and shared capacity. A favourable EBU does not establish feasibility.

\textbf{Field.} The declared state, reference geometry and relevant action relationships used for valuation. Its gradient supplies local directional information. The word does not assert a new fundamental physical force.

\textbf{Field attribution, $R_a$.} A child's signed derivative integral along a declared group path. It sums to the group field reduction under the stated conditions. It is not automatically a uniquely observed causal contribution or an ownership share.

\textbf{Forcing.} A declared external physical change preceding the action baseline in the prospective study. Raw random proposals and physically applied forcing can differ if feasibility rules reject some proposals.

\textbf{Functional balance.} A description of conditions under which relevant functions can continue. A single numerical reference is one representation of part of that idea, not a complete definition for every living system.

\textbf{Gaussian geometry.} The quadratic relative negative-log shape used here to measure standardized deviation. It does not establish an observed normal distribution of actual states, nor an autonomous restoring process.

\textbf{Gradient, $\mu=\nabla V$.} The vector of derivatives of the potential. Each entry describes local change in potential per small change of its coordinate. The gradient is not an actor, a physical stock or a selection algorithm.

\textbf{Group, $G$.} Accepted actions evaluated through a common physical transition because of their declared dependencies. The child actions remain identifiable. A group total does not by itself determine every child allocation.

\textbf{Hessian, $H$.} The matrix of second derivatives. For the separable Gaussian model, $H_{ii}=1/\sigma_i^2$ and off-diagonal entries vanish. Shared action coordinates can still produce group cross terms.

\textbf{Homeostasis.} Maintaining relevant function through continuing change and regulation. The term does not mean that every component is static, that resources are unlimited or that any proposed accounting rule has already achieved regulation.

\textbf{Incidence column, $S_e$.} A signed template turning an edge quantity into its state increment. Multiplying $S_e$ by a finite quantity gives the source and destination changes in stock units.

\textbf{Locality.} A statement about necessary support or information. Mathematical locality, distributed computation and decentralized governance are distinct properties. None should be inferred from the others without a specific argument.

\textbf{Marginal field, $f_e$.} For a lossless directed transfer, $\mu_i-\mu_j$. It gives the initial rate of field reduction per added action quantity. Finite value also includes curvature and any declared process burden.

\textbf{No action.} A genuine zero physical increment with zero process burden. Under the shared-baseline convention it has zero field value. Missing evidence about an action is not the same as verified no action.

\textbf{Potential, $V$.} The declared scalar function of state. In this model, $V=\frac12\sum_i((x_i-x_i^*)/\sigma_i)^2$. It measures represented standardized discrepancy, not total welfare or conserved matter.

\textbf{Process burden, $C$.} Declared physical action burden not already charged through the evaluated state and potential. Its carrier, units and nonduplication must be explained. The principal ideal examples set it to zero.

\textbf{Quantity, $q$.} The finite source-side amount in this edition. A rate $j$ has stock-per-time units and becomes quantity after integration over time. For a constant rate, $q=\Delta t\,j$.

\textbf{Receipt.} A provenance-linked record of a physical event, its accepted comparison, evidence and signed calculation. A record can exist before an institution has decided whether its number changes a transferable balance.

\textbf{Reference, $x_i^*$.} A declared local comparison value. It may be informed by functional evidence, but it is not automatically an observed mean, ethical optimum or dynamic equilibrium.

\textbf{Scale, $\sigma_i>0$.} A fixed positive stock-unit width determining how much a given deviation contributes to the potential. It is separate from measurement uncertainty unless an explicit model establishes a relationship.

\textbf{Service.} A separately specified outcome such as delivery of a required quantity before a deadline. A positive field difference can coexist with unmet service; a negative one can coexist with necessary service.

\textbf{State, $x$.} The ordered collection of represented physical coordinates. Complete state identity includes units, meanings and ordering. Equal potential does not imply equal state.

\textbf{Support, $W$.} Coordinates directly changed by the action. With coupled factors, the set of coordinates required to evaluate the effect can be larger than this write support.

\textbf{Telescoping.} Cancellation of intermediate values in a sum of consecutive differences. It requires consistent definitions and correctly linked boundaries; it does not reconstruct missing physical history.

\textbf{Vector account.} Separate resource-specific results retained together without an undeclared conversion. Even dimensionless scores can have different calibrations and meanings.

\section*{The compact symbol card}
\begingroup\small
\begin{longtable}{P{0.16\linewidth}P{0.38\linewidth}P{0.34\linewidth}}\toprule
Symbol & Object & Units or interpretation\\\midrule
\endfirsthead
\toprule Symbol & Object & Units or interpretation\\\midrule
\endhead
$x_i,z_i$ & Stock at a location & Stock, $\X$\\
$x_i^*$ & Declared reference & Stock, $\X$\\
$\sigma_i$ & Positive reference scale & Stock, $\X$\\
$q$ & Finite action quantity & Stock, $\X$\\
$j$ & Action rate & Stock per time\\
$\delta$ & State increment & Stock in each coordinate\\
$V$ & Gaussian potential & Dimensionless account value\\
$\mu_i$ & Marginal potential & Inverse stock\\
$H$ & Curvature matrix & Inverse squared stock\\
$E$ & Signed finite value & Compatible potential units\\
$R_a$ & Common-path attribution & Compatible potential units\\
$C$ & Residual process burden & Compatible potential units\\
$B_i$ & Proposed action balance & Declared capacity-account units\\
$J$ & External potential audit & Compatible potential units\\\bottomrule
\end{longtable}
\endgroup

Notation alone establishes neither physical calibration nor spending rights.
